
\begin{zusammenfassung}

Diese Arbeit untersucht die Integration von LLMs in robotische Steuerungssysteme zur autonomen Aufgabenplanung und -ausführung. Ein lokal implementiertes LLM dient als logisches Kernsystem, das natürlichsprachliche Anweisungen in strukturierte Befehle übersetzt, welche von einem Roboterarm und analytischen Instrumenten innerhalb eines selbstfahrenden Labors ausgeführt werden. Die vorgeschlagene Architektur schafft eine direkte semantische Schnittstelle zwischen menschlicher Intention und robotischer Aktion und beseitigt dabei herkömmliche Programmierabhängigkeiten. Das Framework wurde in simulierten und realen Experimenten implementiert und validiert. Dabei zeigt es konsistente Interpretation, deterministische Ausführung und geringe Latenzzeiten. Die Ergebnisse bestätigen, dass sprachbasiertes Schlussfolgern eine skalierbare Grundlage für adaptive, menschenverständlich interpretierbare und domänenunabhängige robotische Autonomie bieten kann.

\end{zusammenfassung}
